
\subsubsection*{3.3.1. Corollary}
\phantomsection
\label{sga3:VIIB:3.3.1}

{\itshape Let \(S\) be a locally Noetherian scheme over
\(\mathbb Q\), and let \(G\) be an \(S\)-group scheme that is flat
and locally of finite type.
If \(\omega_{G/S}\) is a locally free \(\mathcal O_S\)-module,
then \(G\) is smooth over \(S\) at every point of the unit
section.}

\textit{Proof.}
Let \(x\) be a point of the unit section, let \(s\) be its image
in \(S\), and denote the local rings at \(x\) and \(s\) by
\(\mathcal O_x\) and \(\mathcal O_s\), respectively.\footnote{%
\label{sga3:VIIB:nde:130}The editors supplied the details of the
argument that follows.}
By hypothesis, \(\mathcal O_s\) and \(\mathcal O_x\) are
Noetherian local rings. The maximal-ideal-adic topology on each
of them therefore agrees with the ``pseudocompact'' topology
defined by the ideals of finite colength. Write
\[
  \widehat{\mathcal O}_x
  \qquad\text{and}\qquad
  \widehat{\mathcal O}_s
\]
for their completions in this topology. By EGA~IV\textsubscript{4},
17.5.3, \(G\) is smooth over \(S\) at \(x\) if and only if
\(\widehat{\mathcal O}_x\) is formally smooth over
\(\widehat{\mathcal O}_s\), both algebras being endowed with their
maximal-ideal-adic topologies. It is therefore enough to prove this
last assertion.

The rings \(\widehat{\mathcal O}_x\) and
\(\widehat{\mathcal O}_s\) are the local rings at \(x\) and \(s\)
of the formal varieties \(\widehat G/\widehat S\) and
\(\widehat S\) defined in \SGARefLink{sga3:VIIB:1.2.6}{1.2.6}. Put
\[
  k=\widehat{\mathcal O}_s,
  \qquad
  H=\operatorname{Spf}\bigl(\widehat{\mathcal O}_x\bigr).
\]
Then \(H\) is an infinitesimal formal \(k\)-variety. Since the
formation of \(\widehat G/\widehat S\) commutes with products
\textup{(loc.\ cit.)}, \(H\) is in fact an infinitesimal formal
\(k\)-group.

Let \(I\) be the augmentation ideal of \(\mathcal O_x\).
By hypothesis,
\[
  \omega_{G/S,x}=I/I^2
\]
is a free \(\mathcal O_s\)-module of finite rank \(n\). Since
\(\mathcal O_s\) is Noetherian, the module \(\omega_{H/k}\),
which is the completion of \(\omega_{G/S,x}\), is identified with
\[
  \omega_{G/S,x}
  \otimes_{\mathcal O_s}
  \widehat{\mathcal O}_s.
\]
It is therefore a topologically free \(k\)-module of rank \(n\).
By the implication \textup{(iv)}\(\Rightarrow\)\textup{\SGARefLink{sga3:vii:paragraph:5-1:equation:ii}{(ii)}} of
\SGARefLink{sga3:VIIB:3.3:cartier-theorem}{Theorem~3.3},
\[
  \widehat{\mathcal O}_x
  \simeq
  k[[\omega_{H/k}]]
  \simeq
  k[[t_1,\ldots,t_n]].
\]
Finally, the latter formal power-series algebra is formally smooth
over \(k\), by EGA~0\textsubscript{IV}, 19.3.3. This proves the
corollary.

Thus we recover a result obtained independently for group schemes
locally of finite presentation over an arbitrary scheme \(S\);
cf.\ VI\textsubscript{B}, \SGARefLink{prop:sga3-VIB-1.6}{1.6}.
